\section{Data Cleaning Detail}\label{app:data}

\Cref{fig:dataoverview} summarizes the cleaned return panel: the distribution of daily returns
pooled across tickers, the cross-sectional correlation structure, the number of tickers
reporting on each session, and the per-ticker volatility spread. The dominant single factor
visible in the correlation panel --- most pairs of VN100 constituents are strongly positively
correlated --- is the structural feature invoked in \cref{sec:insample} to explain why the
robust term concentrates rather than diversifies on this universe.

\begin{figure}[htbp]
\centering
\includegraphics[width=\textwidth]{fig1_data_overview.png}
\caption{Overview of the cleaned VN100 return panel: 98 tickers, 965 sessions,
2022-07-27 to 2026-07-24.}
\label{fig:dataoverview}
\end{figure}

\paragraph{Zero-return diagnostic.}
Beyond the two vendor defects described in \cref{sec:data}, the per-ticker proportion of
exactly-zero return sessions was ranked across all 98 tickers. The purpose was diagnostic: a
single ticker dominating the zero counts would indicate a per-ticker feed problem --- a stock
suspended from trading, or a symbol whose price series is stale --- which is a different defect
requiring a different remedy from the market-wide repeated-price artifact that the $90\%$
filter targets. No such dominance was found; zero returns are distributed across tickers at
levels consistent with ordinary low-liquidity sessions and price-limit days on an exchange with
daily price bands.

\paragraph{Why the holiday filter is statistical.}
A \texttt{dayofweek >= 5} test removes weekends but cannot remove Vietnamese public holidays
that fall on a weekday. Hard-coding a holiday calendar would require annual maintenance and
would fail silently once out of date. The implemented rule --- drop any session on which at
least $90\%$ of tickers post exactly zero return --- keys on the observable signature of the
defect rather than on the calendar. The threshold is deliberately high: genuine market-wide
flat days are possible in principle, and a lower threshold risks discarding real sessions. As
noted in \cref{sec:data}, this filter removes no rows from the present panel.

\section{Additional In-Sample Figures}\label{app:robust}

\begin{figure}[htbp]
\centering
\includegraphics[width=0.6\textwidth]{fig5_robust_effect.png}
\caption{Active positions and Herfindahl--Hirschman index as $\kappa$ increases from $0$ to $1$
with $\gamma$ and $\lambda$ held fixed. Both rise: the portfolio becomes more concentrated, not
more diversified. The mechanism --- degree-one homogeneity of the robust term --- is discussed in
\cref{sec:insample}.}
\label{fig:robust}
\end{figure}

\begin{figure}[htbp]
\centering
\includegraphics[width=0.75\textwidth]{fig4_weights_bar.png}
\caption{Portfolio weights at a small $\lambda$ (denser) against a large $\lambda$ (sparser). The
large-$\lambda$ portfolio concentrates in the few tickers with the best marginal contribution to
the objective; the remaining weights are below the $10^{-4}$ active threshold, not merely small
relative to the active ones.}
\label{app:weights}
\end{figure}

\section{Explicit Objectives for Each Backtest Variant}\label{app:variants}

\Cref{tab:variants} summarizes the six variants of \cref{sec:backtest}; this appendix gives each
one's objective explicitly, each obtained from \cref{eq:main} by fixing or dropping specific
terms --- no variant introduces a term not already present in the general model. A and B
additionally impose $w\ge0$ here, and pick up the concentration constraint of \cref{sec:cap} only
in \cref{tab:beforeafter}. C and D solve no optimization problem at all.

\textbf{A} (robust, long-only) --- $\kappa$ and $\gamma$ searched, $\lambda\equiv0$ (vacuous under
the simplex, \cref{eq:leverage}):
\[
\min_{w\ge0,\ \mathbf1^\top w=1}\ -\muhat^\top w
  \;+\; \underbrace{\kappa\norm{\Ssqrt w}_2}_{\text{robustness}}
  \;+\; \underbrace{\gamma\, w^\top\Sig w}_{\text{variance}}
\]

\textbf{B} (classical, long-only) --- $\kappa\equiv0$ (robust term off), $\gamma$ searched,
$\lambda\equiv0$ (same vacuity as A: B is exactly A with the robust term switched off, nothing
else differs):
\[
\min_{w\ge0,\ \mathbf1^\top w=1}\ -\muhat^\top w
  \;+\; \underbrace{\gamma\, w^\top\Sig w}_{\text{variance}}
\]

\textbf{C} (equal-weight $1/N$) and \textbf{D} (VN100 buy\&hold) --- $w_i = 1/N$ for every $i$, no
objective solved: C resets to $1/N$ every month, D sets it once and lets it drift untouched.

\textbf{E} (sparse-only, long-short) --- $\kappa\equiv0$ (robust off), $\gamma$ searched, $\lambda$
searched over a real grid, the only variant where sparsity is a \emph{designed} choice rather than
a byproduct:
\[
\min_{\mathbf1^\top w=1}\ -\muhat^\top w
  \;+\; \underbrace{\gamma\, w^\top\Sig w}_{\text{variance}}
  \;+\; \underbrace{\lambda\norm{w}_1}_{\text{sparsity / leverage}}
\]

\textbf{F} (full equation, long-short) --- $\kappa$, $\gamma$, $\lambda$ all searched:
\cref{eq:main} exactly as stated, no term fixed and no exposure constraint added, which
\cref{sec:ef} shows is precisely the combination that fails:
\[
\min_{\mathbf1^\top w=1}\ -\muhat^\top w
  \;+\; \underbrace{\kappa\norm{\Ssqrt w}_2}_{\text{robustness}}
  \;+\; \underbrace{\gamma\, w^\top\Sig w}_{\text{variance}}
  \;+\; \underbrace{\lambda\norm{w}_1}_{\text{sparsity / leverage}}
\]

\section{Transaction-Cost Worked Example}\label{app:timeline}

The rolling-window mechanism itself, with the month-numbering table, is given where it is first
used, in \cref{sec:backtest}'s Design subsection (\cref{tab:timeline}). This appendix works
through the turnover and transaction-cost calculation in detail.

At each rebalance, \texttt{src/backtest.py} computes turnover and cost from the full weight
vector as
\begin{lstlisting}[language=Python, numbers=none]
turnover = float(np.abs(w_start - prev).sum())   # sum of |delta w_i| over ALL 98 tickers
cost     = fee * turnover                        # one multiplication
\end{lstlisting}
where \texttt{w\_start} is the new target weight vector for this period and \texttt{prev} is
the \emph{drifted} portfolio at the close of the previous period --- the weights actually held
just before rebalancing, after a month of price movement, not the previous period's target.
This distinction matters: a portfolio that started last month exactly at its target can drift
several points away from it through price movement alone, and correcting that drift is real
turnover the investor must pay for.

\textbf{Example.} Suppose two of the 98 tickers have drifted to $60\%$ and $40\%$ of portfolio
value by the close of the previous period, and the new target for this period is $30\%$ and
$70\%$ (every other ticker unchanged, contributing $0$ to the sum):
\begin{equation}
\text{turnover} = |30\%-60\%| + |70\%-40\%| = 30\% + 30\% = 60\%,
\qquad
\text{cost} = 0.20\% \times 60\% = 0.12\%
\end{equation}
of portfolio value, deducted from the return of the \emph{first trading day} of the new holding
period only --- the portfolio then drifts untouched, with no further trading, for the rest of
the month. For the very first period of the backtest there is no previous portfolio to compare
against, so \texttt{prev} is treated as all zero and turnover $= \sum_i w_i = 1$ exactly (a full
purchase from cash), matching the $100\%$ first-period turnover stated in \cref{sec:backtest}.

\section{Why Transaction Cost Is Kept External to the Objective}\label{app:turnover-design}

\Cref{sec:interpretation} treats transaction cost as a friction applied \emph{after} $w^\star$ is
chosen --- \cref{eq:main} has no term that sees $w_{\text{prev}}$ at all, so the optimizer that
picks $\gamma=1$ this month has no way to know that doing so will force a $60\%$ turnover against
last month's holdings. The natural alternative is to add a turnover penalty
$\eta\norm{w-w_{\text{prev}}}_1$ directly to \cref{eq:main} at each rebalance, letting the
optimizer trade expected return against \emph{both} risk and rebalancing cost simultaneously,
the standard treatment in the transaction-cost portfolio literature. Three reasons this report
keeps cost external instead:

First, \cref{eq:main} is stated and solved once in \cref{sec:insample} as a single-period
problem, and $w_{\text{prev}}$ is meaningful only in the multi-period walk-forward setting of
\cref{sec:backtest} --- folding it in would mean the core model of \cref{sec:modeling} is no
longer the same model actually deployed, which this report's structure deliberately avoids.
Second, $\eta$ would be a fourth hyperparameter selected on the same six-month validation window
already shown (\cref{sec:results7}, \cref{fig:selectedparams}) to select $\gamma$ unstably ---
adding a cost-aware term without first fixing that instability risks fitting $\eta$ to six months
of noise rather than to a genuine cost--return trade-off. Third, and most concretely, the
concentration cap of \cref{sec:cap} was added \emph{because} of a diversification concern
(\cref{app:backtest}), not a cost concern, yet it happens to cut A's and B's turnover as a side
effect --- keeping cost external keeps that attribution clean: \cref{tab:beforeafter}'s
improvement is demonstrably from the cap, not confounded with a simultaneously-tuned cost
penalty. A cost-aware variant, $\eta\norm{w-w_{\text{prev}}}_1$ added once the $\gamma$-selection
instability above is addressed, is left as the natural next extension rather than folded in here.

\section{Backtest Diagnostics}\label{app:backtest}

\begin{table}[htbp]
\centering
\caption{Method B (uncapped), active positions by month --- the trace behind the concentration
constraint of \cref{sec:cap}. Mean $7.56$, range $1$--$14$.}
\label{tab:bmonthly}
\small
\begin{tabular}{@{}lr@{\hspace{18pt}}lr@{\hspace{18pt}}lr@{}}
\toprule
Month & $n$ & Month & $n$ & Month & $n$ \\
\midrule
2024-07 & 12 & 2025-02 & 1  & 2025-09 & 14 \\
2024-08 & 10 & 2025-03 & 2  & 2025-10 & 2  \\
2024-09 & 9  & 2025-04 & 2  & 2025-11 & 9  \\
2024-10 & 9  & 2025-05 & 3  & 2025-12 & 3  \\
2024-11 & 7  & 2025-06 & 12 & 2026-01 & 10 \\
2024-12 & 12 & 2025-07 & 10 & 2026-02 & 11 \\
2025-01 & 12 & 2025-08 & 12 & 2026-03 & 4  \\
        &    &         &    & 2026-04 & 2  \\
        &    &         &    & 2026-05 & 7  \\
        &    &         &    & 2026-06 & 6  \\
        &    &         &    & 2026-07 & 8  \\
\bottomrule
\end{tabular}
\end{table}

\begin{figure}[htbp]
\centering
\includegraphics[width=0.82\textwidth]{fig9_selected_params.png}
\caption{Method A: $(\kappa,\gamma)$ selected by validation Sharpe at each monthly rebalance, and
the resulting turnover. $\gamma$ moves across $\{1, 5, 10\}$ between adjacent periods without
settling, and each move rewrites a meaningful share of the portfolio even under the concentration
cap of \cref{sec:backtest} --- the visual evidence for the turnover diagnosis in
\cref{sec:results7}. $\kappa = 0$ is selected in 18 of 25 periods, consistent with the finding
that the robust term does not earn its keep on this data.}
\label{fig:selectedparams}
\end{figure}

\section{Code Listings}\label{app:code}

The three modules below implement the core contribution of this report: the estimators feeding
\cref{eq:main}, the proximal-subgradient solver of \cref{sec:algorithm}, and the walk-forward
backtest of \cref{sec:backtest}. Full source, with tests, is in the accompanying code submission;
these are reproduced verbatim.

\subsection{\texttt{src/estimators.py}}
\lstinputlisting[language=Python]{code/estimators.py}

\subsection{\texttt{src/prox\_solver.py}}
\lstinputlisting[language=Python]{code/prox_solver.py}

\subsection{\texttt{src/backtest.py}}
\lstinputlisting[language=Python]{code/backtest.py}

\section{Reproduction}\label{app:repro}

\begin{lstlisting}[language={}, basicstyle=\ttfamily\small, numbers=none, frame=single]
python3 -m venv .venv && source .venv/bin/activate
pip install -r requirements.txt
python -m src.data_loader     # only if data/*.parquet is absent
pytest tests/ -v              # 51 passed
python -m nbconvert --to notebook --execute notebook_en.ipynb \
       --output notebook_en.ipynb
\end{lstlisting}

The test suite comprises \textbf{51 tests, all passing}, covering the estimators (symmetry,
positive semidefiniteness, square-root reconstruction, shrinkage bounds), the CVXPY solve
(feasibility, sparsity behavior, the closed-form mean--variance limiting case), and the
backtest (absence of look-ahead in the window construction, turnover accounting, benchmark
correctness). The notebook executes end to end in roughly 50 seconds from the cached parquet
files, with no network access.

The first data download is the only step requiring network access and takes considerably
longer, on the order of tens of minutes, because the free tier of the data API is rate limited
and each of the 98 tickers is fetched separately. Results are cached to
\texttt{data/prices.parquet} and \texttt{data/returns.parquet}, so it is run once.
